% Part: incompleteness
% Chapter: introduction
% Section: overview

\documentclass[../../../include/open-logic-section]{subfiles}

\begin{document}

\olfileid{inc}{int}{ovr}

\olsection{అసంపూర్ణతా ఫలితాల అవలోకనం}

గణితాన్ని ఒక \tetoken{స్వీకృతీకరించదగిన}{!!{axiomatizable}} సిద్ధాంతంలో
సాంకేతికీకరించవచ్చనీ, ఆ సిద్ధాంతం \tetoken{సంపూర్ణమైనదనీ}{!!{complete}},
\tetoken{నిర్ణయించదగినదనీ}{!!{decidable}} నిరూపించవచ్చనీ హిల్బర్ట్ ఆశించాడు.
పైగా, అంతఃప్రజ్ఞావాదపు సవాళ్ల నుండి సాంప్రదాయిక గణితాన్ని రక్షించే అత్యంత
బలహీనమైన ``పరిమితీయ'' పద్ధతులతో ఆ సిద్ధాంతపు అవైరుధ్యాన్ని నిరూపించాలని అతను
లక్ష్యంగా పెట్టుకున్నాడు. ఈ లక్ష్యాలు సాధ్యం కావని గ్యోడెల్ అసంపూర్ణతా సిద్ధాంతాలు
చూపాయి.

రసెల్, వైట్‌హెడ్‌ల \emph{ప్రిన్సిపియా మాథమాటికా}లోని ఒక రూపం
\tetoken{సంపూర్ణం}{!!{complete}} కాదని గ్యోడెల్ మొదటి అసంపూర్ణతా సిద్ధాంతం
చూపింది. కానీ ఆ నిరూపణ నిజానికి అత్యంత సాధారణమైనది; అనేక రకాల సిద్ధాంతాలకు
వర్తిస్తుంది. అంటే \emph{ప్రిన్సిపియా మాథమాటికా} గణితాన్ని సంపూర్ణంగా పట్టుకోలేక
పోయిందని మాత్రమే కాదు; ఆమోదయోగ్యమైన \emph{ఏ} సిద్ధాంతమూ పట్టుకోలేదని అర్థం.
అసంపూర్ణతా సిద్ధాంతాలు వర్తించడానికి సరిపోయే సిద్ధాంత లక్షణాలను వేరుచేయడానికీ,
గ్యోడెల్ నిరూపణ ఈ లక్షణాలపైనే ఆధారపడేలా సాధారణీకరించడానికీ కొంతకాలం పట్టింది.
కానీ ఇప్పుడు అంకగణిత భాష~$\Lang L_A$లోని సిద్ధాంతాలకు మొదటి అసంపూర్ణతా
సిద్ధాంతపు అత్యంత సాధారణ రూపాన్ని చెప్పగలం.
\sourcecorrection{OLTEINCINT-009}{ఆంగ్ల మూలంలోని ``to generalize G\"odel's proof to apply make it depend''లో పరస్పరం కలిసిపోయిన రెండు క్రియాబంధాలున్నాయి. సందర్భానికి అనుగుణంగా, నిరూపణను ఈ లక్షణాలపైనే ఆధారపడేలా సాధారణీకరించడమనే ఉద్దేశిత వాక్యాన్ని ఇచ్చాం.}

\begin{thm}
$\Gamma$ అనేది~$\Lang L_A$లో అవైరుధ్యమైన,
\tetoken{స్వీకృతీకరించదగిన}{!!{axiomatizable}} సిద్ధాంతమై, అన్ని గణనీయ
ప్రమేయాలకూ \tetoken{నిర్ణయించదగిన}{!!{decidable}} సంబంధాలకూ
\tetoken{ప్రాతినిధ్యం వహిస్తే}{!!{represents}}, అప్పుడు $\Gamma$
\tetoken{సంపూర్ణం}{!!{complete}} కాదు.
\end{thm}

$\Gamma$ \tetoken{సంపూర్ణం}{!!{complete}} కాదనడం అంటే, కనీసం ఒక
\tetoken{వాక్యం}{!!{sentence}}~$!A$కు $\Gamma \Proves/ !A$, $\Gamma \Proves/
\lnot !A$ రెండూ కావడం. అటువంటి \tetoken{వాక్యాన్ని}{!!a{sentence}}
($\Gamma)$ నుండి \emph{స్వతంత్రం} అంటాం. స్వతంత్ర వాక్యాలు తప్పకుండా ఉంటాయని
వాస్తవానికి చాలా త్వరగానే నిరూపించవచ్చు. కానీ ఈ సిద్ధాంతానికి గ్యోడెల్ ఇచ్చిన
నిరూపణ శక్తి, అటువంటి స్వతంత్ర \tetoken{వాక్యానికి}{!!{sentence}} ఒక
\emph{నిర్దిష్ట ఉదాహరణను} చూపడంలో ఉంది. ఈ ఆసక్తికర నిర్మాణం~$\Gamma$కు
\emph{గ్యోడెల్ వాక్యం} అనబడే \tetoken{ఒక వాక్యం}{!!a{sentence}}~$!G_\Gamma$ను
ఉత్పత్తి చేస్తుంది. $\Gamma$లో, $!G_\Gamma$ అనేది $!G_\Gamma$ను~$\Gamma$లో
నిరూపించలేమనే వాదనకు తుల్యమైనది కనుక అది నిరూపించలేనిది. ఈ నిర్మాణం
\emph{నిర్మాణాత్మకంగా} అలా చేస్తుంది: $\Gamma$ స్వీకృతీకరణనూ
\tetoken{వ్యుత్పత్తి}{!!{derivation}} వ్యవస్థ వర్ణననూ ఇస్తే, $!G_\Gamma$ను
నిజంగా వ్రాసే పద్ధతిని నిరూపణ ఇస్తుంది.

గ్యోడెల్ నిరూపణలోని నిర్మాణానికి, $\Lang L_A$లోనే $\Lang L_A$ పదాలూ
\tetoken{సూత్రాల}{!!{formula}s} లక్షణాలను, వాటిపై క్రియలను వ్యక్తం చేసే మార్గం
కావాలి. ``$!A$ \tetoken{ఒక వాక్యం}{!!a{sentence}},'' ``$\delta$ అనేది~$!A$కు
\tetoken{ఒక వ్యుత్పత్తి}{!!a{derivation}},'' వంటి లక్షణాలు, $\Subst{!A}{t}{x}$
వంటి క్రియలు వీటిలో ఉన్నాయి. ఈ మార్గం (a) సంకేతాలనూ వాటి క్రమాలనూ (పదాలు,
\tetoken{సూత్రాలు}{!!{formula}s}, \tetoken{వ్యుత్పత్తులు}{!!{derivation}s}
మొదలైనవి ఇవే) సహజ సంఖ్యలుగా ($\Lang L_A$ మాట్లాడగలిగేది వీటి గురించే) చేసే
``సంకేతీకరణ'' ద్వారా ఈ లక్షణాలు, సంబంధాలను వ్యక్తం చేయాలి. (b) సంబంధిత వాస్తవాలను
$\Gamma$ నిరూపించే విధంగా అలా చేయాలి. అందుకే ఈ లక్షణాలు సహజ సంఖ్యల
\tetoken{నిర్ణయించదగిన}{!!{decidable}} లక్షణాలచే సంకేతీకరించబడతాయనీ, క్రియలు
సహజ సంఖ్యలపై గణనీయ ప్రమేయాలకు అనుగుణంగా ఉంటాయనీ చూపాలి. దీన్ని
``వాక్యనిర్మాణపు అంకగణితీకరణ'' అంటారు.

అయితే వాక్యనిర్మాణాన్ని ఎలా అంకగణితీకరించవచ్చో పరిశీలించే ముందు, $\Gamma$
``తగినంత బలమైనది'' అనే షరతును---అంటే అన్ని గణనీయ ప్రమేయాలకూ
\tetoken{నిర్ణయించదగిన}{!!{decidable}} సంబంధాలకూ
\tetoken{ప్రాతినిధ్యం వహిస్తుందనే}{!!{represents}} షరతును---పరిశీలిస్తాం.
దీనికి ``గణనీయం'' అనే పదానికి ఖచ్చిత నిర్వచనం ఇవ్వాలి. ట్యూరింగ్ యంత్రాల నమూనా
ద్వారా గానీ, ఏదైనా సాధారణ-ప్రయోజన ప్రోగ్రామింగ్ భాషలో ప్రోగ్రాములు గణించగల
ప్రమేయాలుగా గానీ, దీన్ని అనేక విధాలుగా చేయవచ్చు. అయితే మన లక్ష్యం ఈ ప్రమేయాలకూ
సంబంధాలకూ $\Lang L_A$ భాషలోని ఒక సిద్ధాంతంలో
\tetoken{ప్రాతినిధ్యం వహించడమే}{!!{represents}s} కనుక, కేవలం సంఖ్యాత్మక
ప్రమేయాల గణనీయతకు సరళమైన నిర్వచనాన్ని ఎంచుకోవడం ఉత్తమం. అదే
\emph{పునరావృత్త ప్రమేయం} అనే భావన. కాబట్టి ముందుగా పునరావృత్త ప్రమేయాలను
చర్చిస్తాం. తరువాత అన్ని పునరావృత్త ప్రమేయాలకూ సంబంధాలకూ ఇప్పటికే $\Th{Q}$
\tetoken{ప్రాతినిధ్యం వహిస్తుందని}{!!{represents}} చూపుతాం. అప్పుడు $\Th{Q}$,
$\Th{PA}$ వంటి నిర్దిష్ట
సిద్ధాంతాలు ``తగినంత బలమైన'' సిద్ధాంతాల ఉదాహరణలని స్థాపించి ఉంటాం కనుక,
వాటికి అసంపూర్ణతా సిద్ధాంతాన్ని వర్తింపజేయగలం.
\sourcecorrection{OLTEINCINT-010}{ఆంగ్ల మూలంలోని ``aim is to represents+s''లో ``to'' తరువాత తప్పుడు క్రియారూపం ఉంది. ముడి OpenLogic గుర్తింపు కీని లక్ష్య టోకెన్‌లో యథాతథంగా ఉంచి, సందర్భానికి సరైన ``ప్రాతినిధ్యం వహించడం'' అనే అర్థాన్ని ఇచ్చాం.}

వాక్యనిర్మాణపు అంకగణితీకరణ చివరికి, \tetoken{సూత్రాలను}{!!{formula}s} సంఖ్యలుగా
సంకేతీకరించడం ద్వారా~$\Gamma$ స్వీకృతాల నుండి నిరూప్యతను వ్యక్తం చేసే
\tetoken{ఒక సూత్రం}{!!a{formula}}~$\OProv[\Gamma](x)$ను ఇస్తుంది. నిర్దిష్టంగా,
$!A$ను సంఖ్య~$n$ సంకేతీకరిస్తే, $\Gamma\Proves !A$ అయినప్పుడు
$\Gamma\Proves\Prov[\Gamma](\num{n})$. $\Gamma$కు ఈ ``నిరూప్యతా విధేయం,''
$\Gamma$ అవైరుధ్యాన్ని కూడా ఒక అర్థంలో~$\Lang L_A$లోని
\tetoken{ఒక వాక్యంగా}{!!a{sentence}} వ్యక్తం చేయడానికి వీలు కల్పిస్తుంది:
$\Gamma$కు ``అవైరుధ్య వాక్యం''గా \tetoken{వాక్యం}{!!{sentence}}
$\lnot\Prov[\Gamma](\num{n})$ను తీసుకుందాం; ఇక్కడ $n$ అనేది ఒక వైరుధ్యపు,
ఉదాహరణకు~$\lfalse$ యొక్క సంకేతం. అవైరుధ్యమైన
\tetoken{స్వీకృతీకరించదగిన}{!!{axiomatizable}} సిద్ధాంతాలు తమ స్వంత అవైరుధ్య
వాక్యాలను కూడా నిరూపించవని రెండవ అసంపూర్ణతా సిద్ధాంతం చెబుతుంది. ఈ సిద్ధాంతం
వర్తించడానికి కావలసిన షరతులు, సిద్ధాంతం అన్ని గణనీయ ప్రమేయాలకూ
\tetoken{నిర్ణయించదగిన}{!!{decidable}} సంబంధాలకూ ప్రాతినిధ్యం వహించడంకంటే
కొంచెం కఠినమైనవి; కానీ $\Th{PA}$ వాటిని నెరవేరుస్తుందని చూపుతాం.

\end{document}
